\uuid{ngKV}
\titre{Famille libre/Génératrice/Base}
\auteur{hanani ; mzari}
\organisation{exo7}
\licence{CC BY-NC-SA 4.0}
\datecreate{2020-01-20}
\niveau{L1}
\module{Algèbre}
\chapitre{Espace vectoriel}
\sousChapitre{Base}
\difficulte{1}
\theme{Base et dimension, Famille libre, Génératrice, Base}
\exo7qcmid{lille-2.espaces-vectoriels.026}

\contenu{
\begin{question}
\qclassification{106.04}
\qtags{motcle=Famille libre/Génératrice/Base}
Dans $\Rr^3$, on considère les vecteurs $u_1=(1,1,0), u_2=(0,1,-1)$ et $ u_3=(-1,0,-1)$. Quelles sont les assertions vraies ?
\begin{answers}
\bad{$\{u_1,u_2,u_3\}$ est une famille libre.}
\bad{$\{u_1,u_2,u_3\}$ est une famille génératrice de $\Rr^3$.}
\good{$u_3$ est une combinaison linéaire de $u_1$ et $u_2$.}
\bad{$\{u_1,u_2,u_3\}$ est une base de $\Rr^3$.}
\end{answers}
\begin{explanations} On vérifie que $u_3= u_2-u_1$, donc $\{u_1,u_2,u_3\}$ n'est pas libre. Par conséquent, $\{u_1,u_2,u_3\}$
n'est pas génératrice de $\Rr^3$, sinon, $\{u_1,u_2\}$ serait aussi génératrice de $\Rr^3$, ce qui contredirait
le fait que toute famille génératrice de $\Rr^3$ doit contenir au moins $3$ vecteurs non nuls.
\end{explanations}
\end{question}
}
